% §1 Introduction
Voice-based conversational agents are no longer a laboratory curiosity: they answer
telephone calls, operate bank accounts, and route service requests on a daily basis
\citep{openai2024realtimeapi,google2024geminivoice}.
Two architectural families dominate this space.
\emph{Native voice models} process speech end-to-end inside a single network, eliminating
inter-component boundaries and their latencies \citep{defossez2024moshi}.
\emph{Cascade architectures} decompose the agent into speech recognition, a text-based
language model, and speech synthesis, exposing transcripts and structured outputs that make
tool use, debugging, and governance practical \citep{yao2023react,langgraph2024}.
Modern cascades are valued precisely because they compose commodity components, swap
models independently, and call external tooling and databases mid-conversation.

The price of the cascade is latency.
Perceptual studies converge on a turn-taking norm of roughly $200$~ms
\citep{levinson2015timing,stivers2009universals}: responses slower than this are perceived as
delayed.
Our companion paper \citep{latencyfloor2025} formalises the \emph{onset} side of responsiveness for cascades: the
Latency Floor Theorem states that the silence gap from end of user speech to first audible
agent audio has an irreducible lower bound, and the Cascade Competitiveness Criterion states
when a cascade can meet the $200$~ms norm.
That work leaves the pipeline's middle stage as a black box and, crucially, addresses only
\textbf{when the agent starts to speak}.

Responsiveness has a second, equally visible half: \textbf{when the agent can fully stop}.
A user who interrupts a banking agent to correct an account number experiences, in that
instant, how much sound the agent keeps making and how long before the system is again
listening.
On tool-heavy turns this stopping behaviour dominates the felt quality of the assistant.
Yet the cessation side is neither standardly defined nor formally bounded: the turn-taking
literature treats barge-in operationally \citep{raux2008optimizing,skantze2021turn}, and
engineers measure barge-in behaviour ad hoc without a shared definition or a lower bound.

This paper supplies the missing floor.
We introduce a two-floor model of realtime responsiveness, prove a phase-conditioned
\emph{cessation floor} for cascade agents, and show how tool round-trips couple to the onset
floor unless they are scheduled onto a background execution clock.
We prove an API-serializability lemma that makes \emph{tool persistence under cancellation}
safe: an interrupted turn cancels generation while pending tool batches commit atomically.
We argue that decoupled-clocks (event-driven) scheduling is the idiomatic configuration that
owns both floors — a sufficiency claim, not a uniqueness claim.
We validate the required semantics with a small, tested, event-driven orchestrator and anchor
the quantitative statements with a monotonic-clock smoke measurement.

Our contributions are:
\begin{enumerate}
  \item a two-floor model of responsiveness (onset floor from our companion work
        \citep{latencyfloor2025}, cessation floor introduced here), together with
        Proposition~\ref{prop:floor_dependency} stating the dependency structure between the
        floors;
  \item a phase-conditioned cessation floor (generate / playback / tool-wait) with a
        detect--decide--cease--resume decomposition and separated \emph{can stop}
        (physical) and \emph{should stop} (policy) statements
        (Theorem~\ref{thm:cessation_floor});
  \item a tool-hiding coupling result showing that tool round-trips sit on the onset side
        unless dispatched to a background execution clock (Theorem~\ref{thm:tool_hiding});
  \item an API-serializability lemma making tool persistence under cancellation safe
        (Lemma~\ref{lem:serializability}) and a filler-placement corollary
        (Corollary~\ref{cor:filler});
  \item a small, tested, event-driven orchestrator providing faithful semantics, plus a
        measured smoke anchor of the cessation floor.
\end{enumerate}

The rest of the paper is structured as follows.
Section~\ref{sec:background} positions the work against incremental dialogue processing,
endpointing, interruption, and orchestration literature.
Section~\ref{sec:framework} defines the two-floor framework and the event system.
Section~\ref{sec:analysis} states and proves the results.
Section~\ref{sec:practical} answers the practitioner's question of which cascade
architecture to build.
Section~\ref{sec:open_problems} lists open directions.
Section~\ref{sec:discussion} discusses scope and limitations, and
Section~\ref{sec:conclusion} concludes.
